% !TEX root = ../main.tex
\begin{abstract}
Sequence saturation can leave a branch in a phylogeny without enough signal to justify the split it represents. SatuTe tests this by estimating a coherence coefficient for the second largest eigenvalue of the substitution model. We ask whether, under rate-heterogeneous Tree-of-Life models, the remaining non-stationary directions should contribute to the same test, and how an explicit invariant-site component changes the infinite-branch null. The relative eigenvalue-weighted statistic keeps the SatuTe null hypothesis and reduces to SatuTe under the Jukes--Cantor model. For \(+I\) mixtures, the invariant category remains jointly distributed across an infinite focal branch and must enter only the null denominator. In a dense paired \(+I+G4\)/\(+I+R4\) grid of 10.8 million simulated alignments, the refined invariant-null statistic gave a mean nominal false-positive rate of 4.92\%, compared with 90.0\% for the counterfactual construction. After independent null calibration, the largest power gains were 11.6 percentage points for \(+I+G4\) and 16.3 percentage points for \(+I+R4\), localized to the detection boundary. The refinement therefore corrects a structural \(+I\) null mismatch and can modestly improve boundary power; it does not change \(+G\)-only analyses.
\end{abstract}
